
This paper asks: when should we retire a benchmark? Among 115 Papers With Code benchmarks, a statistically significant structural break in residual CoV occurs at paper\_count\* $\approx$ 39 --- the point at which benchmark competition dynamics shift from high-variance exploration to low-variance saturation. Before this threshold, performance scores vary 3.$81\times$ more than the global baseline; after it, variance falls to 0.$20\times$ of the pre-breakpoint level. This 80\% variance collapse (measured relative to the pre-breakpoint segment) is confirmed by two independent statistical tests (permutation p=0.035, piecewise F-test p=0.0021) and corroborated by three experiments with distinct test designs and null hypotheses.

Contributions: (1) first detected structural break in PwC benchmark residual CoV at paper\_count\*$\approx 39$ via PELT change-point detection; (2) empirical characterization of both regimes (exploration: 3.$81\times$ variance, saturation: 0.$20\times$ compression), each confirmed at p < 0.01; (3) a demonstrated methodology --- OLS detrend + PELT + permutation test --- applicable to any leaderboard dataset for saturation onset detection.

Future directions grounded in the experimental evidence include: threshold sensitivity analysis across min\_papers values, longitudinal within-benchmark CoV trajectories for causal attribution, domain-stratified paper\_count\* analysis, and generalization to HELM and OpenLLM leaderboard data.

